\begin{otherlanguage*}{english}
\noindent\textbf{Chapter XVI, manuscript pages 1--124: equivalences before models.}

Grothendieck returns to the axioms because he wants to understand which part of a homotopy theory belongs to the equivalences themselves. A useful factorization supplies a way of calculating a homotopy fibre or a homotopy pushout; it should not determine what that construction means. This conviction explains both the long inventory in these pages and its recurrent changes of emphasis. The object being studied is a class $W$ of functors between small categories. From it one wants to recover the localization $\mathrm{Hot}_W$, categories of coherent coefficients over every index category, and the operations relating those categories. The chapter asks how much geometry is forced by a few closure properties of $W$, and how much must still be supplied by an axiom of descent, a factorization, or a smallness condition.

The letter to Thomason of 2 April 1991 states this position particularly clearly: a category with weak equivalences is the underlying datum, while cofibrations and fibrations are additional means of constructing its homotopy theory. The same letter recalls the observation that a derived category ``ne vient jamais seule'': diagrams of complexes have their own derived categories, whose relationships carry information lost on passing immediately to diagrams in a single derived category \citep[\S\S1--2]{Grothendieck1991Letter}. Chapter XVI belongs to that project. Its geometric language of fibres, open and closed subcategories, and integration is a language for constructing coherent diagram theories. The connection with the theory of coefficients is structural, rather than an analogy added after the constructions have been made.

\medskip\noindent\textbf{The first distinction: absolute and relative localization.}

The starting axioms are deliberately modest. They impose the specified saturation conditions, declare a category with a terminal object to be aspherical, and require an aspherical functor to be an equivalence. Here asphericity of $u:A\to B$ is tested on the comma categories $A/b$. The stronger axiom Loc\,(3 bis) concerns a map $A\to B$ over an arbitrary third category $S$, and tests the maps $A/s\to B/s$. This difference governs much of the chapter. The first condition controls individual objects and functors with contractible fibres in the comma sense. The second permits the comparison of two varying families. In later terminology the first system defines a \emph{weak} fundamental localizer; the usual definition of a fundamental localizer incorporates the relative condition.

The comparison with split Grothendieck fibrations in \S1 makes this distinction geometric. Replacing a category over $S$ by its free split fibration turns comma categories into actual fibres. The comparison with the original category has an adjoint and is therefore already harmless for the weak axioms. What remains is exactly the question whether fibrewise equivalences of two families are equivalences of their total categories. The commentary spells out the variance in this replacement: Grothendieck fibrations are described by contravariant category-valued functors, and their fibres are compared with undercategories over the base. Interchanging these conventions silently would change the argument.

The connectedness axiom of \S2 rules out the degenerate theories in which all nonempty categories become indistinguishable. It allows $\pi_0$ to descend through localization and explains the special role of discrete categories. The converse characterization by nontriviality invokes an earlier product argument whose editors explicitly require (3 bis); its weak-axiom generality is therefore not established by that reference. This section also contains a useful warning about localization arguments. Adjunctions that preserve the designated equivalences descend by their units and counits. An identification between the localization of an infinite product of relative categories and the product of their localizations requires additional justification. The finite and infinite cases must therefore be separated. Moreover, the literal formulation of Lemma 2.5 with an arbitrary class of arrows omits an identity hypothesis needed even in the finite case. The audit isolates this defect without disturbing the intended application to localizers, which already contain the identities.

\medskip\noindent\textbf{Excision and the two meanings of fibration.}

The next two sections study complementary ways of making constructions invariant. For pushouts, the intrinsic class $\mathrm{Cof}_W$ consists of maps along which cobase change preserves $W$-equivalences. The forward Mayer--Vietoris condition is an excision principle for a cover by two sieves, or dually by two cosieves. Dwyer inclusions supply enough explicit maps with which to work. Brown's lemma is the decisive bridge from stability of trivial cofibrations to stability of arbitrary equivalences under these pushouts \citep{Brown1973}. The commentary gives the missing argument, retaining the distinction between Grothendieck's sketch, the cited general result, and its application in $\mathrm{Cat}$.

Forward excision must be distinguished from the reverse implication, which tries to recover an equivalence on the intersection from equivalences on the two pieces and their union. The latter is substantially stronger. It fails for ordinary weak homotopy equivalences, and the blanket assertion in Theorem 5.4 that the ordinary localizer satisfies Loc\,(6) must be read with this defect exposed: its forward part is valid, its reverse part is not. Cone and suspension constructions explain why this is a mathematical issue rather than a matter of terminology. They also give a useful further development: reverse cube excision forces the inversion of integral homology equivalences. The detailed proof appears at the end of the commentary on \S3, with its precise scope and the literature checked.

For pullbacks, $\mathrm{Fib}_W$ is defined intrinsically by preservation of $W$-equivalences under base change. It is not the class of fibrations in an arbitrarily chosen Quillen model structure. Nor is it the class of Grothendieck fibrations with $W$-equivalent transition functors. The latter is an explicit source of candidate maps, and \S4 asks when those candidates belong to the intrinsic class. Universal equivalences, weak fibrations, perfect maps, and the auxiliary localizer $W_0\subseteq W$ provide progressively different tests. The many implications labelled A, B, C, D and F record which tests are sufficient for which calculations. Their hypotheses matter: a statement about maps over a fixed base does not automatically assert stability under every change of target, and the existence of enough intrinsic fibrations does not identify a prescribed geometric class with those fibrations.

These pages are exploratory, but their main mechanism is robust. A factorization by a sufficiently well behaved class reduces preservation of all equivalences to preservation of trivial maps in that class. Repeated path constructions then turn local constancy into the stronger comparison needed for arbitrary base change. The commentary follows this mechanism through the diagram of implications, provides the missing intermediate verifications, and distinguishes genuine proof gaps from incorrectly copied indices and from slips already present in the edition. A worked analysis of component monodromy gives a concrete test for the difference between locally constant Grothendieck fibrations and intrinsic fibrations for $\pi_0$-equivalences.

\medskip\noindent\textbf{What the minimality theorem changes.}

The most consequential later result for these pages is stronger than the minimality theorem in its familiar formulation. Cisinski proves that ordinary weak homotopy equivalences form the least \emph{weak} fundamental localizer \citep[Theorem 6.1.18, pp.~242--243]{Cisinski2006}. Consequently Grothendieck's $W_\omega$, defined here using only the weak axioms, equals $\mathbf W_\infty$. This settles the suggestions that $W_\omega$ might fail the filtered-colimit axiom, and that a fundamental localizer strictly contained in $\mathbf W_\infty$ might retain simple homotopy information. It is essential to cite this stronger theorem: the relative theorem alone would not justify its application to the definition used at the beginning of Chapter XVI.

Theorem 6.1 gives a second consequence. Its fibration hypothesis, together with connectedness, forces $W\subseteq\mathbf W_\infty$. Minimality supplies the opposite inclusion. The other branch can be proved directly: a disconnected $W$-aspherical category, after perfect replacement of a chosen point, produces an empty fibre; universal base change then forces every functor into $W$. This supplies the missing dependency without invoking the general nontriviality criterion. Consequently $D_{W_\omega}$ permits exactly the ordinary localizer and the total localizer; Loc\,(5), or non-totality under this fibration hypothesis, selects the former. The eight conditions in the main rectangle are therefore equivalent, while F remains strictly weaker. This explains why the stronger fibre calculations exclude the truncated and homological examples that otherwise belong naturally to the theory. The finite truncation localizers do have enough intrinsic fibrations: their proper generalized Thomason models answer the question about F positively \citep[\S9.2.1, Corollary 9.2.15 and Theorem 5.2.15]{Cisinski2006}.

One qualification is necessary for the list on ms.~87: its bare Loc\,(6), if read as including reverse excision, is incompatible with the simultaneous connectedness and fibration requirements. Equality with $\mathbf W_\infty$ describes the intended forward-excision version, not that inconsistent literal list.

The finite-poset calculation in \S6 nevertheless retains its independent interest. Ordinary homotopy of realizations is weaker than homotopy through natural transformations of functors. Subdivision preserves the realization but need not provide a categorical homotopy inverse. The commentary proves the cycle example explicitly and distinguishes categorical homotopy, weak homotopy type, and simple homotopy type. Minimality prevents the proposed refinement from being encoded by a fundamental localizer on all small categories; it does not identify these three notions or erase the finer finite combinatorics.

\medskip\noindent\textbf{Smallness and the passage to derivators.}

Accessibility in \S5 supplies the finiteness principle needed to construct models. The relevant smallness is a bounded-subcategory property: a prescribed small part of an equivalence can be enlarged, without exceeding the chosen cardinal bound, to a small equivalence. The argument for $\mathbf W_\infty$ uses a countable succession of homotopy witnesses. The notation $\pi_0$ in the cardinal estimate is expressly used for $\aleph_0$ and is not a transcription error. Translating this passage into modern presentability language requires care about successor cardinals; one must not identify a countable-category bound with finite presentability. The commentary gives an explicit closure construction and supplies the cofinal-tail choice needed to represent a basepoint in a filtered system.

The final section returns to the goal of the axioms. It asks when $\mathrm{HOT}_W$ has homotopy Kan extensions, detects isomorphisms pointwise, and satisfies the required exactness properties. The source's numbering of the derivator axioms is kept separate from modern numbering. In particular its decomposition axiom concerns finite sums. Its assertion for arbitrary relative categories is therefore not refuted by a failure involving infinite products. Likewise, the calculation of a colimit at an initial object in \S7.7 is correct with the manuscript's contravariant diagram convention.

Two later theorems give firm landmarks. Accessible fundamental localizers in the relative sense admit the generalized Thomason model structures constructed by Cisinski; left properness is part of the theorem, whereas right properness requires the corresponding properness condition \citep[Theorem 5.2.15]{Cisinski2006}. Suitable model categories produce derivators, with homotopy Kan extensions and their pointwise formulas \citep{Cisinski2003,Groth2013}. These results answer the model-existence problem under explicit hypotheses. They do not validate the circular inference in Corollary 7.7.2, where the identification of $W$ with the localizer detected by $\mathrm{HOT}_W$ uses the very relative axiom it is meant to establish. The edition preserves the warning about that circle, and so does the audit. For the remaining weak-axiom variants, the results found are kept apart from the implications that the sources checked do not settle.

\end{otherlanguage*}
